%-------------------------------------------------------------
% Subhadeep Koley -- Research CV (1-2 pages)
% Compile with: pdflatex Subhadeep_Koley_CV.tex
%-------------------------------------------------------------
\documentclass[10pt,letterpaper]{article}

\usepackage[margin=0.7in]{geometry}
\usepackage[T1]{fontenc}
\usepackage{charter}
\usepackage{titlesec}
\usepackage{enumitem}
\usepackage{xcolor}
\usepackage[hidelinks]{hyperref}
\usepackage{tabularx}

\definecolor{accent}{HTML}{1F4E79}
\hypersetup{colorlinks=true, urlcolor=accent, linkcolor=accent}

\pagestyle{empty}
\setlength{\parindent}{0pt}

% Section headings: small caps with a rule underneath
\titleformat{\section}{\large\scshape\color{accent}}{}{0em}{}[\vspace{-6pt}\rule{\textwidth}{0.4pt}]
\titlespacing*{\section}{0pt}{10pt}{6pt}

% Tight bullet lists
\setlist[itemize]{leftmargin=1.2em, itemsep=1pt, topsep=2pt, parsep=0pt}

% \entry{Title}{Dates}{Subtitle}{Location}
\newcommand{\entry}[4]{%
  \textbf{#1} \hfill #2 \\
  \textit{#3} \hfill \textit{#4}\par\vspace{1pt}}

% \pub{Authors}{Title}{Venue}{Links}
\newcommand{\pub}[4]{%
  \item #1. \textbf{#2}. \textit{#3}. #4}

\newcommand{\me}{\textbf{Subhadeep Koley}}

\begin{document}

%----------------------------- HEADER -----------------------------
\begin{center}
  {\LARGE \textbf{Subhadeep Koley}}\\[4pt]
  PhD Student, Computer Science and Engineering, Lehigh University\\[2pt]
  \href{mailto:svk324@lehigh.edu}{svk324@lehigh.edu} \;|\;
  \href{https://subhadeepk.github.io}{subhadeepk.github.io} \;|\;
  \href{https://scholar.google.com/citations?user=_b0J6xoAAAAJ&hl=en}{Google Scholar} \;|\;
  \href{https://github.com/subhadeepk}{github.com/subhadeepk} \;|\;
  \href{https://www.linkedin.com/in/subhadeep-koley-70251b1bb/}{LinkedIn}
\end{center}

\vspace{-4pt}
\textbf{Research interests:} aerial robotics, physical human-robot interaction, motion planning,
learning-based control. Seeking research internships for Summer 2027.

%----------------------------- EDUCATION -----------------------------
\section{Education}

\entry{Lehigh University}{Aug 2024 -- Present}
      {PhD, Computer Science and Engineering. Advisor: Prof.\ David Salda\~na}{Bethlehem, PA}
\vspace{3pt}
\entry{Indian Institute of Engineering Science and Technology (IIEST), Shibpur}{2020 -- 2024}
      {B.Tech., Mechanical Engineering. CGPA: 8.56/10}{Howrah, India}

%----------------------------- PUBLICATIONS -----------------------------
\section{Publications}

\textit{Under review}
\begin{itemize}
  \pub{\me, Benjamin Greenberg, Yifei Simon Shao, Juan Aceros, Nadia Figueroa, David Salda\~na}
      {HINT-Blimp: Human INTent Inference from Multimodal Cues for Robotic Blimps}
      {Under review, ICRA 2027}
      {[\href{https://arxiv.org/abs/2609.27154}{arXiv}] [\href{https://youtu.be/b_QbHGdgBf4}{video}]}
  \pub{\me, Subhrajit Bhattacharya, David Salda\~na}
      {Planning Trajectories that Bounce: Reflection Classes for Collision-Tolerant Robots}
      {Under review, ICRA 2027}
      {[\href{https://arxiv.org/abs/2609.27145}{arXiv}] [\href{https://youtu.be/bygwxGIbFh4}{video}]}
\end{itemize}

\textit{Conference papers}
\begin{itemize}
  \pub{Yizhuo Wang, Yuhong Cao, Jimmy Chiun, \me, Mandy Pham, Guillaume Adrien Sartoretti}
      {ViPER: Visibility-based Pursuit-Evasion via Reinforcement Learning}
      {Conference on Robot Learning (CoRL), 2024}
      {[\href{https://openreview.net/forum?id=EPujQZWemk}{paper}] [\href{https://www.youtube.com/watch?v=r2SurmnPKuA}{video}] [\href{https://github.com/marmotlab/ViPER}{code}]}
  \pub{Ajit Das, Saumendra Nath Mishra, \me, Sourav Sarkar, Achintya Mukhopadhyay, Swarnendu Sen}
      {Early Detection of Thermal Runaway of Lithium-ion Battery: An Experimental Study}
      {IEEE Conf.\ on Sustainable Energy and Future Electric Transportation (SEFET), 2023}
      {\textbf{Best Paper Award}. [\href{https://doi.org/10.1109/SEFET57834.2023.10245991}{paper}]}
\end{itemize}

\textit{Workshop papers}
\begin{itemize}
  \pub{\me, Benjamin Greenberg, Yifei Simon Shao, Nadia Figueroa, David Salda\~na}
      {HINT-Blimp: Human INTent Inference from Multimodal Cues for Robotic Blimps}
      {IROS 2026 Workshop on Nonverbal Cues for Human-Robot Cooperative Intelligence, Pittsburgh, 2026.
       Spotlight talk and poster}
      {[\href{https://nocworkshop.github.io/2026-IROS/resources/15/Camera_ready_HintBlimp_IROS_NoC_workshop.pdf}{paper}]
       [\href{https://www.youtube.com/watch?v=P_ZU_f_-qek}{video}]
       [\href{https://nocworkshop.github.io/2026-IROS/resources/15/IROS_workshop_poster.pdf}{poster}]}
\end{itemize}

\textit{Posters and demos}
\begin{itemize}
  % TODO: replace with the exact title of your NERC poster
  \item \textbf{HINT-Blimp}: Poster presentation, \textit{Northeast Robotics Colloquium (NERC)},
        Princeton University, Oct 2026.
  \item Live demonstration of \href{https://ieeexplore.ieee.org/document/11127424}{MochiSwarm}, an open-source swarm of lightweight robotic blimps. \textit{ICRA 2025}, Atlanta.
\end{itemize}

%----------------------------- RESEARCH EXPERIENCE -----------------------------
\section{Research Experience}

\entry{Graduate Researcher, Swarms Lab, Lehigh University}{Aug 2024 -- Present}
      {Advisor: Prof.\ David Salda\~na}{Bethlehem, PA}
\begin{itemize}
  \item Working on autonomous navigation and mapping of large unknown indoor environments using Unitree Go2.
  \item Built a multimodal intent-inference system for robotic blimps using physical pushes, voice and pointing;
        combining pushes and voice identified the intended goal 86\% of the time, usually within two interactions,
        and let users steer the blimp around obstacles only they knew about.
  \item Developed a planner that lets collision-tolerant robots bounce off walls deliberately, grouping paths into
        reflection classes by wall-contact sequence; the best bouncing path reduced control effort by a median
        16.6\% versus a collision-free path in simulation.
  \item Won first place at \href{https://engineering.lehigh.edu/news/article/lehigh-robotics-team-soars-victory-defend-republic-drone-competition}{Defend the Republic 2024} against teams including Georgia Tech, UMich and ASU;
        built and maintained 16+ robots and added ultrasonic wall avoidance to the hardware and software stack.
\end{itemize}

\entry{Research Intern, MARMoT Lab, National University of Singapore}{Jan 2024 -- May 2024}
      {Mentor: Dr.\ Guillaume Sartoretti}{Remote}
\begin{itemize}
  \item Worked on reinforcement learning for multi-robot visibility-based pursuit-evasion in unknown
        environments; published at CoRL 2024 (ViPER).
\end{itemize}

\entry{SURGE Research Intern, SPiN Lab, IIT Kanpur}{May 2023 -- Aug 2023}
      {Mentor: Dr.\ Ketan Rajawat}{Kanpur, India}
\begin{itemize}
  \item Part of a 10-person team building ROS software for a decentralized drone swarm, deployed on five real drones
        [\href{https://drive.google.com/drive/folders/1TChMz4PDQW4giv97GUWhrWedk10NNIpQ}{report}]
        [\href{https://youtu.be/Jn5sE0hCgE8}{video}].
  \item Developed a mission framework supporting leader-follower flight and independent per-drone missions
        [\href{https://youtu.be/xq_rb1A7OfA}{video}].
  \item Built decentralized leader election so the swarm elects a new leader on leader failure.
  \item Developed inter-drone collision avoidance, tested under motion capture, and network monitoring and
        congestion control for the swarm's mesh network.
\end{itemize}

\entry{Research Intern, Jadavpur University}{Jul 2022 -- Dec 2022}
      {Mentor: Dr.\ Sourav Sarkar}{Kolkata, India}
\begin{itemize}
  \item Built a heat- and gas-sensor setup for early detection of lithium-ion battery thermal runaway,
        leading to a Best Paper Award at IEEE SEFET 2023.
  \item Built a multi-point wireless temperature acquisition system for a Tata Steel-funded project
        studying cooling rates in rolling-mill cooling beds.
\end{itemize}

\entry{Flight Software Intern, EndureAir Systems}{Dec 2022 -- Feb 2023}
      {Mentor: Dr.\ Kuldeep Kumar Dhiman}{Noida, India}
\begin{itemize}
  \item Worked with the PX4 flight stack; built a MATLAB range-extension tool for a drop tank with delayed wing deployment.
\end{itemize}

%----------------------------- AWARDS -----------------------------
\section{Awards}
\begin{itemize}
  \item \href{https://engineering.lehigh.edu/news/article/lehigh-robotics-team-soars-victory-defend-republic-drone-competition}{First place}, Defend the Republic drone competition, Indiana University \hfill Nov 2024
  \item Best Paper Award, IEEE SEFET \hfill Aug 2023
\end{itemize}

%----------------------------- TEACHING -----------------------------
\section{Teaching}
\begin{itemize}
  \item Teaching Assistant, CSE 340: Design and Analysis of Algorithms, Lehigh University
\end{itemize}

%----------------------------- SKILLS -----------------------------
\section{Technical Skills}
% TODO: add C++, ROS 2, or other languages/tools you actually use in the blimp work.
\begin{tabularx}{\textwidth}{@{}l X@{}}
  \textbf{Languages}   & Python, MATLAB \\
  \textbf{Robotics}    & ROS, PX4, Gazebo, CoppeliaSim, MoveIt, motion capture \\
  \textbf{ML / Vision} & PyTorch, OpenCV \\
  \textbf{Hardware}    & ESP32, ATmega328, Raspberry Pi, SolidWorks, 3D printing \\
\end{tabularx}

\end{document}